%Please consider the following section of the "Formvorschriften für die gedruckte Version"
%Im Anhang ist eine Zusammenfassung (Abstract) mitzubinden.
%Ist die Arbeit in einer Fremdsprache verfasst, ist im Anhang jedenfalls eine deutsche Zusammenfassung mitzubinden.
\chapter{Abstract}
The kinematics of the Taurus molecular cloud complex have so far mainly been studied through the
discrete velocities of its individual members. A full three-dimensional space velocity, however,
requires a radial velocity, which is available for only about one third of the known members, while
proper motions are measured for essentially all of them. This thesis therefore reconstructs a
continuous three-dimensional velocity field of the complex from the full stellar sample, using
information field theory as implemented in \texttt{NIFTy}. Each velocity component is modeled as a
correlated field on a grid of $3~\mathrm{pc}$ cells spanning $126 \times 69 \times 81~\mathrm{pc}$,
with a power spectrum inferred from the data. The stars enter the
likelihood through their observables, so that a star without a radial velocity still constrains the
field within the plane of the sky while its missing component is supplied by the spatial correlation
with its neighbors. The input catalog consists of 557 members identified with \texttt{SigMA} in
Gaia~DR3 astrometry, of which 123 contribute a radial velocity after quality cuts.

The reconstruction reproduces the bulk motion of the complex, resolves velocity differences between
its subgroups and recovers an age--velocity gradient that is already visible in the stellar data.
Both the mean motion and the quantities derived from the field, however, depend strongly on the
individual run of the algorithm: repetitions with a different random seed, with different quality cuts or with a
modified model yield mean motions that scatter by several $\mathrm{km\,s^{-1}}$, while the posterior
uncertainty within a single run stays far below that. A comparison of several such runs is therefore
used to separate features demanded by the data from artifacts of a particular reconstruction. The
residuals between the stars and the field expose the central limitation of the model: the peculiar
motion of a star has nowhere to go but into the field itself. Granting the stars a common intrinsic
velocity dispersion, inferred alongside the field, brings the weighted residuals of the proper
motions from several standard deviations down to below one, although it does not bring the field
closer to the data.
